\chapter{Références et Bibliographie}
\label{chap:references}

Les algorithmes et méthodes d'informatique graphique et de vision par ordinateur implémentés dans cette plateforme s'appuient sur les travaux fondamentaux suivants :

\begin{thebibliography}{9}

\bibitem{levoy1988}
Levoy, Marc. 
\textit{Display of surfaces from volume data.}
IEEE Computer graphics and Applications 8.3 (1988): 29-37.
\newblock (Fondation théorique du Direct Volume Rendering et du Raymarching utilisés dans le \texttt{VolumeViewer}).

\bibitem{kabsch1976}
Kabsch, Wolfgang. 
\textit{A solution for the best rotation to relate two sets of vectors.}
Acta Crystallographica Section A: Crystal Physics, Diffraction, Theoretical and General Crystallography 32.5 (1976): 922-923.
\newblock (Algorithme de stabilisation Rigid-Body par SVD utilisé dans \texttt{registration.py}).

\bibitem{threejs}
Cabello, Ricardo, et al.
\textit{Three.js: A JavaScript 3D library.}
\newblock \url{https://threejs.org/} (Moteur 3D sous-jacent au \texttt{VolumeViewer} et \texttt{TrackingViewer}).

\bibitem{openseadragon}
OpenSeadragon Contributors.
\textit{OpenSeadragon: An open-source, web-based viewer for zoomable images.}
\newblock \url{https://openseadragon.github.io/} (Moteur de rendu pyramidal 2D utilisé par le \texttt{DeepZoomViewer}).

\bibitem{imaris}
Bitplane / Oxford Instruments.
\textit{Imaris: 3D/4D Image Visualization and Analysis Software.}
\newblock \url{https://imaris.oxinst.com/} (Inspiration pour l'UX du mode \textit{Emission/MIP} et la gestion des canaux fluorescents).

\end{thebibliography}
